
% =====================================================================
\section{Non-return and the cessation of inquiry}
\label{sec:inquiry}
% =====================================================================

We close the development with two order-theoretic results. The first
(\Cref{thm:nonreturn}) shows that the record of committed individuations is
monotone under every traversal, so no process returns to a prior individuation
state and resemblance never amounts to identity. The second
(\Cref{thm:cessation}) shows that, because thinning a separator has diverging
cost and bounded return, a finite agent halts at the floor; knowledge thereafter
attaches to bounded wholes and never to sub-floor parts. Neither result uses any
dynamics: both are statements about the order on committed separations.

\subsection{The committed record}

\begin{definition}[Committed record]
\label{def:record}
Fix a refinement chain of realisable partitions
$\Part_0\succeq\Part_1\succeq\cdots$. The \emph{committed record} after a finite
sequence of individuation steps is the count
\[
M\;:=\;\#\{\text{separator cells committed across all steps so far}\}\in\N,
\]
the number of cells moved out of an approximant's separator over the process
\textup{(}\Cref{def:process}, \Cref{cor:monotone-commit}\textup{)}.
\end{definition}

\begin{theorem}[Non-return]
\label{thm:nonreturn}
The committed record $M$ is monotone non-decreasing under every traversal of the
refinement chain. No traversal returns the record to a previously held smaller
value; in particular, a step intended to undo a prior commitment is itself a
committed step and increases $M$. Consequently two individuation states with the
same configuration of cells but different records are distinct, and resemblance
of configuration never entails identity of state.
\end{theorem}

\begin{proof}
By \Cref{cor:monotone-commit} each committing step increases the count of
committed separator cells by at least one, and the refinement order
$\Part_j\succeq\Part_{j+1}$ is one-directional: to ``undo'' a commitment is to
perform a further act---a re-coarsening or re-drawing---which by
\Cref{def:process} is itself a committing step and so increments $M$. Hence $M$
never decreases; any attempted return raises it. If two states $s,s'$ share a
configuration of cells but have records $M(s)\neq M(s')$, then they are
distinguished by the record itself: there is a partition of states, namely
``record $\le M(s)$ versus record $>M(s)$,'' separating them. By the principle
that states are identical only when no realisable partition distinguishes them,
$s\neq s'$. Resemblance of configuration is therefore compatible with difference
of record, and identity requires both.
\end{proof}

\begin{remark}[Resemblance versus identity]
\label{rem:resemblance}
\Cref{thm:nonreturn} is the order-theoretic core of irreversibility for
individuation: a process may revisit a configuration that \emph{looks like} an
earlier one, but the committed record that accompanies it has strictly grown, so
the state is new. The monotone record is the invariant that configuration alone
cannot capture. Nothing here invokes any flow or external parameter; the
monotonicity is a property of the order on commitments.
\end{remark}

\subsection{Cessation}

\begin{theorem}[Inquiry cessation]
\label{thm:cessation}
Let $\Gamma$ be a refinement cost \textup{(}\Cref{def:cost}\textup{)} and let the
return of an individuation---the additional content it correctly certifies---be
bounded by $\mu(\Space)<\infty$ \textup{(}\ref{brs:bdd}\textup{)}. Then for any
agent with a positive threshold on cost-per-return there is a finite thickness
$\thick_\ast>0$ at which further thinning is rejected: the agent halts at a
strictly positive floor and certifies nothing below it. Knowledge so obtained
attaches to bounded wholes of thickness $\ge\thick_\ast$ and never to sub-floor
parts.
\end{theorem}

\begin{proof}
By \Cref{thm:cost}, realising thickness $t$ costs $\Gamma_t\ge g(t)$ with
$g(t)\to+\infty$ as $t\to0^+$. The return is bounded above by $\mu(\Space)$, a
finite constant. Hence the marginal return per unit cost,
$\mu(\Space)/g(t)$, tends to $0$ as $t\to0^+$. An agent that continues thinning
only while cost-per-return stays below a fixed positive threshold $\theta$ stops
at the largest $t$ with $g(t)\le\mu(\Space)/\theta$; call it $\thick_\ast$. Since
$g$ is finite away from $0$, $\thick_\ast>0$. Below $\thick_\ast$ the agent does
not refine, so it certifies no region of thickness $<\thick_\ast$; by
\Cref{thm:detectability} regions of scale below the floor are in any case
engulfed. Thus knowledge attaches to wholes of thickness $\ge\thick_\ast$ and to
no sub-floor part.
\end{proof}

\begin{remark}[Why knowledge attaches to wholes]
\label{rem:wholes}
\Cref{thm:cessation} gives the final shape of the picture. Individuation is by
negation (\Cref{thm:negation-unique}); it is non-instantaneous and leaves a
residue (\Cref{thm:noninstant}); the residue cannot be exhibited or certified
(\Cref{thm:residue-not-measure}); the record of commitments only grows
(\Cref{thm:nonreturn}); and the cost of pursuing the residue to zero diverges
against a bounded return, so a finite agent rationally halts at a positive floor
(\Cref{thm:cessation}). What such an agent can know are bounded wholes---regions
that clear their own boundary thickness---attributed and reasoned about as units.
It cannot know its sub-floor parts, not for want of effort but because below the
floor a part is indistinguishable from its boundary and its certification is
self-referentially obstructed. The boundary of knowledge is the floor, and the
floor is positive.
\end{remark}
